\section*{वर्णमाला --- Phonemic Inventory \& Script Guide}
\vspace{0.2em}
{\small\color{bpgray}The complete phonemic inventory of Bhojpuri in Kaithi script, Devanagari, and Latin transliteration.}
\vspace{0.4em}

\subsection*{\color{bpblue}स्वर आ चिह्न / Vowels \& Signs}

{\bfseries\small Independent Vowels (स्वतंत्र स्वर)}
\renewcommand{\arraystretch}{1.15}
\begin{longtable}{@{} c c l l l @{}}
\toprule
\textbf{Kaithi} & \textbf{Devanagari} & \textbf{Latin} & \textbf{IPA} & \textbf{Name} \\
\midrule
\endfirsthead
\toprule
\textbf{Kaithi} & \textbf{Devanagari} & \textbf{Latin} & \textbf{IPA} & \textbf{Name} \\
\midrule
\endhead
\bottomrule
\endfoot

\K{\ba} & अ & a & /ə/ & a \\
\K{\baa} & आ & aa & /aː/ & aa \\
\K{\bi} & इ & i & /i/ & i \\
\K{\bii} & ई & ii & /iː/ & ii \\
\K{\bu} & उ & u & /u/ & u \\
\K{\buu} & ऊ & uu & /uː/ & uu \\
\K{\be} & ए & e & /eː/ & e \\
\K{\bai} & ऐ & ai & /ɛː/ & ai \\
\K{\bo} & ओ & o & /oː/ & o \\
\K{\bau} & औ & au & /ɔː/ & au \\
\bottomrule
\end{longtable}

\vspace{0.3em}

{\bfseries\small Dependent Vowel Signs (मात्रा) --- Demonstrated with \K{\bka} (ka)}
\begin{longtable}{@{} c c c l l @{}}
\toprule
\textbf{Matra} & \textbf{Devanagari} & \textbf{Kaithi Result} & \textbf{Latin} & \textbf{IPA} \\
\midrule
\endfirsthead
\toprule
\textbf{Matra} & \textbf{Devanagari} & \textbf{Kaithi Result} & \textbf{Latin} & \textbf{IPA} \\
\midrule
\endhead
\bottomrule
\endfoot

-- (inherent) & -- & \K{\bka} & ka & /kə/ \\
\K{\bmaa} & ा & \K{\bka\bmaa} & kaa & /kaː/ \\
\K{\bmi} & ि & \K{\bka\bmi} & ki & /ki/ \\
\K{\bmii} & ी & \K{\bka\bmii} & kii & /kiː/ \\
\K{\bmu} & ु & \K{\bka\bmu} & ku & /ku/ \\
\K{\bmuu} & ू & \K{\bka\bmuu} & kuu & /kuː/ \\
\K{\bme} & े & \K{\bka\bme} & ke & /keː/ \\
\K{\bmai} & ै & \K{\bka\bmai} & kai & /kɛː/ \\
\K{\bmo} & ो & \K{\bka\bmo} & ko & /koː/ \\
\K{\bmau} & ौ & \K{\bka\bmau} & kau & /kɔː/ \\
\bottomrule
\end{longtable}

\vspace{0.3em}

{\bfseries\small Virama \& Nasal Marks}
\begin{longtable}{@{} c c l l @{}}
\toprule
\textbf{Sign} & \textbf{Devanagari} & \textbf{Name} & \textbf{Function} \\
\midrule
\endfirsthead
\toprule
\textbf{Sign} & \textbf{Devanagari} & \textbf{Name} & \textbf{Function} \\
\midrule
\endhead
\bottomrule
\endfoot
\K{\char"25CC\bvir} & ◌् & Virama (Halant) & Suppresses inherent vowel /a/ \\
\K{\banus} & ◌ं & Anusvara & Nasal mark \\
\K{\bchandra} & ◌ँ & Chandrabindu & Vowel nasalization mark \\
\bottomrule
\end{longtable}

\newpage

\subsection*{\color{bpblue}व्यंजन आ संयुक्ताक्षर / Consonants \& Conjuncts}

{\bfseries\small Sparsh Vyanjan --- Stops \& Nasals (स्पर्श व्यंजन)}
\begin{longtable}{@{} l c c c c c @{}}
\toprule
\textbf{Place} & \textbf{Vl Unasp} & \textbf{Vl Asp} & \textbf{Vd Unasp} & \textbf{Vd Asp} & \textbf{Nasal} \\
\midrule
\endfirsthead
\toprule
\textbf{Place} & \textbf{Vl Unasp} & \textbf{Vl Asp} & \textbf{Vd Unasp} & \textbf{Vd Asp} & \textbf{Nasal} \\
\midrule
\endhead
\bottomrule
\endfoot

Velar (क वर्ग) & \K{\bka} (क) & \K{\bkha} (ख) & \K{\bga} (ग) & \K{\bgha} (घ) & \K{\bnga} (ङ) \\
Palatal (च वर्ग) & \K{\bca} (च) & \K{\bcha} (छ) & \K{\bja} (ज) & \K{\bjha} (झ) & \K{\bnja} (ञ) \\
Retroflex (ट वर्ग) & \K{\btta} (ट) & \K{\bttha} (ठ) & \K{\bdda} (ड) & \K{\bddha} (ढ) & \K{\bnna} (ण) \\
Dental (त वर्ग) & \K{\bta} (त) & \K{\btha} (थ) & \K{\bda} (द) & \K{\bdha} (ध) & \K{\bna} (न) \\
Labial (प वर्ग) & \K{\bpa} (प) & \K{\bpha} (फ) & \K{\bba} (ब) & \K{\bbha} (भ) & \K{\bma} (म) \\
\bottomrule
\end{longtable}

\vspace{0.4em}

{\bfseries\small Antasth \& Ushma --- Approximants \& Fricatives (अन्तस्थ आ ऊष्म)}
\begin{longtable}{@{} c c c c c c c c @{}}
\toprule
\textbf{ya (य)} & \textbf{ra (र)} & \textbf{la (ल)} & \textbf{va (व)} & \textbf{sha (श)} & \textbf{ssa (ष)} & \textbf{sa (स)} & \textbf{ha (ह)} \\
\midrule
\K{\bya} & \K{\bra} & \K{\bla} & \K{\bva} & \K{\bsha} & \K{\bssa} & \K{\bsa} & \K{\bha} \\
/j/ & /r/ & /l/ & /v/ & /ʃ/ & /ʂ/ & /s/ & /h/ \\
\bottomrule
\end{longtable}

\vspace{0.4em}

{\bfseries\small संयुक्ताक्षर / Conjuncts} \\
{\small Two or more consonants join via Virama (\K{\char"25CC\bvir}) to form conjunct clusters.}
\begin{longtable}{@{} c c l l @{}}
\toprule
\textbf{Kaithi} & \textbf{Devanagari} & \textbf{Structure} & \textbf{Latin} \\
\midrule
\endfirsthead
\toprule
\textbf{Kaithi} & \textbf{Devanagari} & \textbf{Structure} & \textbf{Latin} \\
\midrule
\endhead
\bottomrule
\endfoot

\K{\bka\bvir\bka} & क्क & ka + virama + ka & kka \\
\K{\bka\bvir\bkha} & क्ख & ka + virama + kha & kkha \\
\K{\bka\bvir\bga} & क्ग & ka + virama + ga & kga \\
\K{\bka\bvir\bta} & क्त & ka + virama + ta & kta \\
\K{\bka\bvir\bpa} & क्प & ka + virama + pa & kpa \\
\K{\bga\bvir\bka} & ग्क & ga + virama + ka & gka \\
\K{\bga\bvir\bma} & ग्म & ga + virama + ma & gma \\
\K{\bda\bvir\bka} & द्य & da + virama + ya & dya \\
\K{\bta\bvir\bra} & त्र & ta + virama + ra & tra \\
\K{\bsa\bvir\bta} & स्त & sa + virama + ta & sta \\
\bottomrule
\end{longtable}
